\chapter{\LaTeX{}基础语法}
	\section{字体字号设置}
	\begin{description}
	\item{字体}\\
	由于本模板使用的是CTEX封装的\texttt{ctexbook}类为基础的，自带了六款字体，以及自定义了一款字体。字体的使用代码为\texttt{\textbackslash heiti\{\heiti{想改变为黑体的内容}\}}，以及\texttt{\textbackslash songti}, \texttt{\textbackslash fangsong}, \texttt{\textbackslash kaishu}, \texttt{\textbackslash lishu}, \texttt{\textbackslash youyuan}, \texttt{\textbackslash hwxingkai}即可。它们的演示如下：\\
	{\songti 这一行是测试\textcolor{red}{宋体}字体的示范}\\
	{\fangsong 这一行是测试\textcolor{red}{仿宋}字体的示范}\\
	{\lishu 这一行是测试\textcolor{red}{隶书}字体的示范}\\
	{\youyuan 这一行是测试\textcolor{red}{幼圆}字体的示范}\\
	{\kaishu 这一行是测试\textcolor{red}{楷书}字体的示范}\\
	{\heiti 这一行是测试\textcolor{red}{黑体}字体的示范}\\
	{\hwxingkai 这一行是测试\textcolor{red}{华文行楷}字体的示范}\\
	{\fzqingke 方正清刻本悦宋体}\\
	如果需要使用本地的一些特殊字体的话，请参见\url{http://zzgthk.iteye.com/blog/994662}来定义。\\
	\item{字号}\\
	而字号的话，则可以用\texttt{\textbackslash zihao\{-4\}}，以及\texttt{\textbackslash zihao\{4\}}来表示小四号字体和四号字体，同理可以从小五号\texttt{\textbackslash zihao\{-5\}}字体到初号\texttt{\textbackslash zihao\{0\}}字体。\\
	\item{间距}\\
	对于间距的控制，\LaTeX{}中的间距比Office要强大得多了！Office除了空格回车外，差不多就没有了。而对于\LaTeX{}来说，可使用的间隔控制命令有：\texttt{\textbackslash quad},\texttt{\textbackslash qquad}，它们分别表示间隔一个汉字和两个汉字距离，以及配合\texttt{\textbackslash hspace\{3cm\}}则表示水平空格3cm距离，同理还有\texttt{\textbackslash vspace\{3cm\}}，以及\texttt{\textbackslash hfill}表示水平填充一个距离可拉伸的空白距离。\LaTeX{}支持的单位有：\texttt{in}---英寸(inch)，\texttt{mm}---毫米，\texttt{cm}---厘米，\texttt{pt}---Points(大约1\slash 72~inch)，以及相对距离：\texttt{em}---当前字体的字符“\textbf{M}”的宽度，\texttt{ex}---当前字体的字符“\textbf{x}”的宽度，长度可以为负值，如~-1.5em~。\\
	另外，\LaTeX{}会忽略掉你所写的论文中的空格和换行符。如果需要空格的地方请用“$\sim$”，需要换行的地方请用“\textbackslash\textbackslash ”来控制，每段段末请加“\texttt{\textbackslash par}”来表示一个段落的结束。换行的段落结束的区别在于换行后不会有缩进，而段落结束的话，下一段会自动缩进两个汉字距离。
	
	\end{description}

	
	\section{章节标题设置}
	\section{图形和多图排列}
	有些时候论文需要将多个图形排列在一起，成2x2或2x3的样子，以下就举例来说明这个用法。
	
\begin{figure}[htb]
	\subfigure[small box with a long caption]{\label{fig:minia}\includegraphics[width=.3\textwidth]{buaamark.eps}}\hfill
	\subfigure[big box]{\label{fig:minib}\includegraphics[width=.3\textwidth]{buaamark.eps}}
	\caption{Mutipicture inside subfigure}
	\label{subfig1}
\end{figure}	
这是另外一种多图的形式。\par

\begin{figure}[htb]
\begin{minipage}{0.5\textwidth}
	\centering
	\includegraphics[width=0.35\textwidth]{buaamark.eps}
	\caption{SubfigureA}	
	\label{minifiga}
\end{minipage}
\begin{minipage}{0.5\textwidth}
	\centering
	\includegraphics[width=0.35\textwidth]{buaamark.eps}
	\caption{SubfigureB}	
	\label{minifigb}
\end{minipage}
\caption{另外一种多图的表现}
\end{figure}
或者排列成矩形的形式，如图\ref{maxfig}所示。\par
\begin{figure}[htb]
\centering
	\begin{tabular}{cc}
	\subfigure[Figure:A]{\label{suba}\includegraphics[width=.3\textwidth]{buaamark.eps}}\qquad & \subfigure[Figure:B]{\label{subB}\includegraphics[width=.3\textwidth]{buaamark.eps}}\\
	\subfigure[Figure:C]{\label{subc}\includegraphics[width=.3\textwidth]{buaamark.eps}}\qquad & \subfigure[Figure:D]{\label{subd}\includegraphics[width=.3\textwidth]{buaamark.eps}}\\
	\end{tabular}
\caption{矩形的subfig排列}
\label{maxfig}

\end{figure}

\clearpage


	\section{表格和长表格}
	如果您要排版的表格长度超过一页，那么推荐使用 \textsf{longtable} 或者 \textsf{supertabular} 
宏包，模板对 \textsf{longtable} 进行了相应的设置，所以用起来可能简单一些。
表~\ref{example:longtable} 就是 \textsf{longtable} 的简单示例。
\begin{center}
\begin{longtable}{|l|l|l|}
\caption{长表格示例} \label{example:longtable} \\

\hline \multicolumn{1}{|c|}{\textbf{Time (s)}} & \multicolumn{1}{c|}{\textbf{Triple chosen}} & \multicolumn{1}{c|}{\textbf{Other feasible triples}} \\ \hline 
\endfirsthead

\multicolumn{3}{l}%
{{\bfseries \tablename\ \thetable{} --接\,上\,页}} \\
\hline \multicolumn{1}{|c|}{\textbf{Time (s)}} &
\multicolumn{1}{c|}{\textbf{Triple chosen}} &
\multicolumn{1}{c|}{\textbf{Other feasible triples}} \\ \hline 
\endhead

\hline \multicolumn{3}{|r|}{{接\,下\,页}} \\ \hline
\endfoot

\hline \hline
\endlastfoot

0 & (1, 11, 13725) & (1, 12, 10980), (1, 13, 8235), (2, 2, 0), (3, 1, 0) \\
2745 & (1, 12, 10980) & (1, 13, 8235), (2, 2, 0), (2, 3, 0), (3, 1, 0) \\
5490 & (1, 12, 13725) & (2, 2, 2745), (2, 3, 0), (3, 1, 0) \\
8235 & (1, 12, 16470) & (1, 13, 13725), (2, 2, 2745), (2, 3, 0), (3, 1, 0) \\
10980 & (1, 12, 16470) & (1, 13, 13725), (2, 2, 2745), (2, 3, 0), (3, 1, 0) \\
13725 & (1, 12, 16470) & (1, 13, 13725), (2, 2, 2745), (2, 3, 0), (3, 1, 0) \\
16470 & (1, 13, 16470) & (2, 2, 2745), (2, 3, 0), (3, 1, 0) \\
19215 & (1, 12, 16470) & (1, 13, 13725), (2, 2, 2745), (2, 3, 0), (3, 1, 0) \\
21960 & (1, 12, 16470) & (1, 13, 13725), (2, 2, 2745), (2, 3, 0), (3, 1, 0) \\
24705 & (1, 12, 16470) & (1, 13, 13725), (2, 2, 2745), (2, 3, 0), (3, 1, 0) \\
27450 & (1, 12, 16470) & (1, 13, 13725), (2, 2, 2745), (2, 3, 0), (3, 1, 0) \\
30195 & (2, 2, 2745) & (2, 3, 0), (3, 1, 0) \\
32940 & (1, 13, 16470) & (2, 2, 2745), (2, 3, 0), (3, 1, 0) \\
35685 & (1, 13, 13725) & (2, 2, 2745), (2, 3, 0), (3, 1, 0) \\
38430 & (1, 13, 10980) & (2, 2, 2745), (2, 3, 0), (3, 1, 0) \\
41175 & (1, 12, 13725) & (1, 13, 10980), (2, 2, 2745), (2, 3, 0), (3, 1, 0) \\
43920 & (1, 13, 10980) & (2, 2, 2745), (2, 3, 0), (3, 1, 0) \\
46665 & (2, 2, 2745) & (2, 3, 0), (3, 1, 0) \\
49410 & (2, 2, 2745) & (2, 3, 0), (3, 1, 0) \\
52155 & (1, 12, 16470) & (1, 13, 13725), (2, 2, 2745), (2, 3, 0), (3, 1, 0) \\
54900 & (1, 13, 13725) & (2, 2, 2745), (2, 3, 0), (3, 1, 0) \\
57645 & (1, 13, 13725) & (2, 2, 2745), (2, 3, 0), (3, 1, 0) \\
60390 & (1, 12, 13725) & (2, 2, 2745), (2, 3, 0), (3, 1, 0) \\
63135 & (1, 13, 16470) & (2, 2, 2745), (2, 3, 0), (3, 1, 0) \\
65880 & (1, 13, 16470) & (2, 2, 2745), (2, 3, 0), (3, 1, 0) \\
120780 & (1, 13, 16470) & (2, 2, 2745), (2, 3, 0), (3, 1, 0) \\
123525 & (1, 13, 13725) & (2, 2, 2745), (2, 3, 0), (3, 1, 0) \\
126270 & (1, 12, 16470) & (1, 13, 13725), (2, 2, 2745), (2, 3, 0), (3, 1, 0) \\
129015 & (2, 2, 2745) & (2, 3, 0), (3, 1, 0) \\
131760 & (2, 2, 2745) & (2, 3, 0), (3, 1, 0) \\
134505 & (1, 13, 16470) & (2, 2, 2745), (2, 3, 0), (3, 1, 0) \\
137250 & (1, 13, 13725) & (2, 2, 2745), (2, 3, 0), (3, 1, 0) \\
139995 & (2, 2, 2745) & (2, 3, 0), (3, 1, 0) \\
142740 & (2, 2, 2745) & (2, 3, 0), (3, 1, 0) \\
145485 & (1, 12, 16470) & (1, 13, 13725), (2, 2, 2745), (2, 3, 0), (3, 1, 0) \\
148230 & (2, 2, 2745) & (2, 3, 0), (3, 1, 0) \\
150975 & (1, 13, 16470) & (2, 2, 2745), (2, 3, 0), (3, 1, 0) \\
153720 & (1, 12, 13725) & (2, 2, 2745), (2, 3, 0), (3, 1, 0) \\
156465 & (1, 13, 13725) & (2, 2, 2745), (2, 3, 0), (3, 1, 0) \\
159210 & (1, 13, 13725) & (2, 2, 2745), (2, 3, 0), (3, 1, 0) \\
161955 & (1, 13, 16470) & (2, 2, 2745), (2, 3, 0), (3, 1, 0) \\
164700 & (1, 13, 13725) & (2, 2, 2745), (2, 3, 0), (3, 1, 0) \\
\end{longtable}
\end{center}
	\section{使用~listing~显示代码}
原先~ctexbook~文档类和~listings~宏包配合使用时，代码在换页时会出现莫名其妙的错误，后来经高人指点，顺利解决了。
感兴趣的话，可以看看\href{http://bbs.ctex.org/viewthread.php?tid=53451}{这里}。
这里给使用~listings~宏包插入源代码的例子，这里是一段C代码。
另外，listings宏包真可谓博大精深，可以实现各种复杂、漂亮的效果，想要进一步学习的同学，可以参考
\href{http://mirror.ctan.org/macros/latex/contrib/listings/listings.pdf}{listings宏包手册\footnote{此节完全参考武汉理工大学论文模板}}。

\begin{lstlisting}[language={C}, caption={一段C源代码}]
#include <stdio.h>
#include <unistd.h>
#include <sys/types.h>
#include <sys/wait.h>
% 这一行是测试中文支持的。
int main() {
  pid_t pid;

  switch ((pid = fork())) {
  case -1:
    printf("fork failed\n");
    break;
  case 0:
    /* child calls exec */
    execl("/bin/ls", "ls", "-l", (char*)0);
    printf("execl failed\n");
    break;
  default:
    /* parent uses wait to suspend execution until child finishes */
    wait((int*)0);
    printf("is completed\n");
    break;
  }

  return 0;
}
\end{lstlisting}

再给一个插入MATLAB代码的例子，感谢~daisying~站友提供的代码。

\begin{lstlisting}[language={matlab}, caption={一段MATLAB源代码}]
function paper1
r=0.05;
n=100;
T=1;
X=1;
v0=0.8;
sigma=sqrt(0.08);
deltat=T/n;
for i=1:n
    t(i)=i*deltat;
    w(i)=random('norm',0,t(i),1);
end
for i=1:n
    alpha(i)=0.39;
end
for i=1:n
    temp=0;
    for k=1:i
        temp=temp+alpha(k);
    end
    B(i)=exp(r*t(i));
    BB(i)=B(i)*exp(temp*deltat);
    BBB(i)=exp(-r*(T-t(i)));
end
for i=1:n
    s0(i)=X*BBB(i);
    v(i)=v0*exp((r-0.5*sigma^2)*t(i)+sigma*w(i));
    for j=i+1:n
        D=X*BBB(j);
        d1=(log(v(i)/D)+(r+sigma^2/2)*(t(j)-t(i)))/(sigma*sqrt(t(j)-t(i)));
        d2=d1-(sigma*sqrt(t(j)-t(i)));
        ppp(i,j)=D*exp(-r*(t(j)-t(i)))*(1-cdf('normal',d2,0,1))-v(i)*(1-cdf('n
ormal',d1,0,1));
    end
end
for i=1:n
    s1(i)=0;
    for j=i+1:n
        s1(i)=s1(i)+BB(j)^(-1)*alpha(j)*deltat*(X*BBB(j)-B(j)/B(i)*ppp(i,j));
    end
    s2(i)=0;
    for j=1:n
        s2(i)=s2(i)+alpha(j);
    end
    s2(i)=X*exp(-r*T-s2(i)*deltat);
    s(i)=BB(i)*(s1(i)+s2(i));
end
plot(s)
hold on;
plot(s0);
\end{lstlisting}